\chapter{Fault Pattern Catalogue}
\label{app:patterns}

This appendix documents the fourteen requirement-related fault patterns (RF01--RF14)
encoded in the HSP-Agile pattern guard.  The catalogue is derived from the
Agile-SOFL fault taxonomy of \citet{li2023iet} and operationalised as
static analysis rules over generated Python implementations.  Each pattern is
task-aware: detectors receive both the candidate source and the FSF task
artefact (signature, scenario list, external variables) so that structural
checks can be grounded in specification context rather than generic code
heuristics alone.

% ---------------------------------------------------------------------------
\section{Pattern Taxonomy Overview}
\label{sec:pat:taxonomy}
% ---------------------------------------------------------------------------

The pattern identifiers RF01--RF14 form a closed taxonomy covering seven
\emph{category} dimensions that reflect the locus of the fault relative to
the specification--implementation boundary:

\begin{description}
  \item[\textbf{requirement}.]
        Misalignment between stated preconditions or documentation and the
        guard structure actually implemented (RF01, RF14).
  \item[\textbf{fsf}.]
        Violations of FSF-specific structural conventions: default-branch
        handling and scenario coverage (RF02, RF09).
  \item[\textbf{interface}.]
        Omission of declared output variables from return values (RF04).
  \item[\textbf{logic}.]
        Incorrect conditional structure, spurious completeness, or vacuous
        branches (RF05, RF07, RF13).
  \item[\textbf{boundary}.]
        Mishandling of zero-valued inputs or off-by-one threshold comparisons
        (RF06, RF11).
  \item[\textbf{implementation}.]
        Coding-quality defects that do not necessarily violate a single guard
        but indicate unreliable translation from specification to code
        (RF03, RF08, RF12).
  \item[\textbf{sofl}.]
        Neglect of external program variables declared in the SOFL process
        signature (RF10).
\end{description}

Each pattern is additionally assigned a \emph{severity} rating on a four-level
ordinal scale:

\begin{itemize}
  \item \textbf{critical} (1 pattern): defects that indicate the implementation
        cannot reflect scenario-dependent behaviour at all; acceptance is
        blocked unconditionally.
  \item \textbf{high} (8 patterns): defects with strong correlation to
        specification non-conformance on boundary or overlap inputs; included
        in the high-severity acceptance set.
  \item \textbf{medium} (4 patterns): quality or traceability concerns that are
        recorded in the audit trace but do not block acceptance.
  \item \textbf{low} (1 pattern): stylistic redundancy unlikely to cause
        functional failure in isolation.
\end{itemize}

The hybrid acceptance predicate of Definition~\ref{def:accept} consults only
patterns whose severity is \emph{high} or \emph{critical}.  We denote this
subset the \emph{high-severity set}:
\begin{equation}
  \mathcal{P}^H \;=\;
    \{\text{RF01},\;\text{RF02},\;\text{RF04},\;\text{RF05},\;\text{RF06},\;
      \text{RF07},\;\text{RF09},\;\text{RF11},\;\text{RF13}\},
  \label{eq:app:high-sev}
\end{equation}
comprising nine patterns.  A candidate $c$ is pattern-safe on task $t$ when
$\sum_{p \in \mathcal{P}^H} \delta_p(c, t) = 0$.  Medium- and low-severity
hits are appended to the structured feedback message for the next synthesis
iteration but do not veto acceptance.  Table~\ref{tab:pat:severity-summary}
summarises the severity partition.

\begin{table}[htbp]
  \caption{Severity partition of the pattern catalogue ($|\mathcal{P}| = 14$).}
  \label{tab:pat:severity-summary}
  \centering
  \thesistable{%
    \small
    \setlength{\tabcolsep}{4pt}
    \begin{tabularx}{\linewidth}{@{}lXc@{}}
      \toprule
      Severity   & Pattern IDs                          & Count \\
      \midrule
      Critical   & RF07                                 & 1 \\
      High       & RF01, RF02, RF04, RF05, RF06, RF09, RF11, RF13 & 8 \\
      Medium     & RF03, RF08, RF10, RF14               & 4 \\
      Low        & RF12                                 & 1 \\
      \bottomrule
    \end{tabularx}%
  }
\end{table}

Detection is implemented as a hybrid of Python abstract syntax tree (AST)
traversal and regular-expression matching over source text.  AST-based checks
are preferred for structural invariants (branch counts, return-key sets,
else-branch presence); regex checks capture surface-form regularities that are
expensive to express as tree predicates (magic numeric literals, empty handlers,
missing docstrings).  When parsing fails, a synthetic \code{SYNTAX} match with
critical severity is emitted and the candidate is rejected immediately.

% ---------------------------------------------------------------------------
\section{Complete Pattern Reference}
\label{sec:pat:reference}
% ---------------------------------------------------------------------------

Table~\ref{tab:pat:full-catalogue} provides the authoritative reference for
all fourteen patterns.  The \emph{Detection Method} column indicates the
primary mechanism employed by the pattern guard; task-context checks (e.g.\
comparing branch count against the number of FSF scenarios) are noted as
\emph{AST + task context}.

\footnotesize
\setlength{\tabcolsep}{3pt}
\begin{longtable}{@{}p{0.055\linewidth}p{0.13\linewidth}p{0.08\linewidth}p{0.065\linewidth}p{0.27\linewidth}p{0.27\linewidth}@{}}
  \caption{Complete HSP-Agile fault pattern catalogue (RF01--RF14).}
  \label{tab:pat:full-catalogue} \\
  \toprule
  \textbf{ID} & \textbf{Name} & \textbf{Category} & \textbf{Sev.} &
  \textbf{Description} & \textbf{Detection Method} \\
  \midrule
  \endfirsthead

  \multicolumn{6}{c}{\tablename\ \thetable{} --- \textit{continued from previous page}} \\
  \toprule
  \textbf{ID} & \textbf{Name} & \textbf{Category} & \textbf{Sev.} &
  \textbf{Description} & \textbf{Detection Method} \\
  \midrule
  \endhead

  \midrule
  \multicolumn{6}{r}{\textit{Continued on next page}} \\
  \endfoot

  \bottomrule
  \endlastfoot

  RF01 &
  Missing precondition check &
  requirement &
  high &
  Input validation is absent before the first assignment or return in the
  function body; guard conjuncts declared in the FSF test predicate are not
  reflected in a leading conditional. &
  Regex: function entry without a leading \code{if} guard; spec hint
  cross-checks FSF test predicates. \\[4pt]

  RF02 &
  Wrong default branch &
  fsf &
  high &
  The implementation lacks an explicit \code{else} branch to handle the FSF
  \texttt{others} scenario, causing fall-through or implicit defaults that
  diverge from the specification. &
  AST: absence of \code{else}/\code{elif} chain terminus when task declares an
  \texttt{others} scenario. \\[4pt]

  RF03 &
  Hardcoded magic number &
  implementation &
  medium &
  A numeric literal appears in a \code{return} statement without derivation
  from specification constants or input variables; excludes canonical
  boolean-like values \code{0} and \code{1}. &
  Regex: \code{return <digit>} with exclusion of \code{return 0/1}. \\[4pt]

  RF04 &
  Missing output field &
  interface &
  high &
  The return dictionary omits one or more output variables declared in the
  process signature, yielding incomplete structured outputs. &
  AST + task context: compare return-dict keys against declared output names. \\[4pt]

  RF05 &
  Swapped comparison &
  logic &
  high &
  A comparison operator in the implementation is suspected to be reversed
  relative to the FSF guard (e.g.\ \code{<} where \code{>} is specified). &
  Regex over assignment/comparison sites; guard extraction from FSF for
  cross-reference (see Section~\ref{subsec:method:detection}). \\[4pt]

  RF06 &
  No guard for zero input &
  boundary &
  high &
  Zero-valued inputs that appear as explicit boundaries in FSF guards
  (\code{eq 0}, \code{gt 0}) are not handled by a dedicated conditional
  branch. &
  Regex: function signature present without zero-check branch; spec hint
  references FSF zero-boundary scenarios. \\[4pt]

  RF07 &
  Unconditional success &
  logic &
  \textbf{critical} &
  The function returns a fixed success indicator (e.g.\
  \code{success = 1}) on all paths without evaluating any guard-dependent
  conditional, indicating spurious completeness. &
  AST: constant success return with no \code{if} nodes in the function body. \\[4pt]

  RF08 &
  Missing type coercion &
  implementation &
  medium &
  Return values in dictionary form are not cast to integer type where the
  signature declares natural-number outputs. &
  Regex: dictionary return without explicit integer cast. \\[4pt]

  RF09 &
  Incomplete FSF coverage &
  fsf &
  high &
  The number of conditional branches in the implementation is strictly less
  than the number of non-default FSF scenarios, indicating omitted scenario
  handlers. &
  AST + task context: branch count $<$ scenario count (adjusted for
  \texttt{others}). \\[4pt]

  RF10 &
  External variable ignored &
  sofl &
  medium &
  External program variables declared in the SOFL process specification are
  never referenced in the implementation body. &
  AST + task context: name-set intersection of \code{ext} declarations and
  referenced identifiers. \\[4pt]

  RF11 &
  Off-by-one boundary &
  boundary &
  high &
  A strict inequality (\code{>}) is used where a weak inequality (\code{>=})
  is required, or vice versa, shifting the effective guard region by one
  boundary point. &
  Regex: \code{> 0} pattern; threshold literals cross-referenced against FSF
  guard constants. \\[4pt]

  RF12 &
  Duplicate assignment &
  implementation &
  low &
  The same output variable is assigned more than once within a single branch,
  with the earlier assignment shadowed and therefore dead. &
  Regex: repeated assignment to the same identifier within a code region. \\[4pt]

  RF13 &
  Empty handler &
  logic &
  high &
  An \code{else} branch contains only \code{pass}, leaving a scenario handler
  vacuous and producing no output assignment. &
  Regex: \code{else: pass} at line end; AST confirmation of empty body. \\[4pt]

  RF14 &
  Spec-comment mismatch &
  requirement &
  medium &
  The function lacks a docstring referencing the FSF process, reducing
  traceability between specification text and implementation and indicating
  possible copy-paste synthesis without spec grounding. &
  Regex: function definition not followed by a triple-quoted docstring. \\

\end{longtable}
\normalsize

% ---------------------------------------------------------------------------
\section{Specification-Confusion Patterns}
\label{sec:pat:sc}
% ---------------------------------------------------------------------------

\emph{Specification confusion} denotes a class of faults in which the
implementer misreads, misorders, or incompletely internalises the FSF
specification while producing syntactically valid code.  Such faults often
yield implementations that appear structurally reasonable---complete conditional
chains, defined outputs on every branch---yet assign incorrect semantics to
guard regions because scenario priority, boundary thresholds, or default
behaviour were interpreted incorrectly.

Seven patterns in the catalogue directly target specification-confusion
failure modes.  Table~\ref{tab:pat:spec-confusion} lists them with their
principal confusion mechanism and the specification-mutation operators under
which they are most frequently observed in the prevention evaluation protocol
(Appendix~\ref{app:benchmark}).

\begin{table}[htbp]
  \caption{Patterns associated with specification-confusion faults.}
  \label{tab:pat:spec-confusion}
  \centering
  \thesistable{%
    \footnotesize
    \setlength{\tabcolsep}{3pt}
    \begin{tabularx}{\linewidth}{@{}p{0.06\linewidth}p{0.20\linewidth}X X@{}}
      \toprule
      \textbf{ID} & \textbf{Confusion mechanism} & \textbf{Typical symptom} &
      \textbf{Related spec mutants} \\
    \midrule
    RF01 &
    Partial guard extraction &
    Only a subset of conjuncts in a compound FSF test predicate are
    implemented &
    SNO (predicate negation) \\
    RF02 &
    Default-scenario misidentification &
    The \texttt{others} clause outputs are assigned to the wrong branch &
    DRO (else removal), MCO (scenario drop) \\
    RF05 &
    Operator inversion &
    Relational operators are reversed relative to the specification &
    ORO (operator swap), ICO (condition inversion) \\
    RF06 &
    Zero-boundary neglect &
    Inputs at zero are not routed to the scenario that explicitly guards
    \code{eq 0} &
    BCO (boundary weakening) \\
    RF09 &
    Scenario omission &
    One or more ordered scenarios are not represented as conditional branches &
    MCO (scenario drop) \\
    RF11 &
    Threshold displacement &
    Strict/weak inequality confusion shifts guard regions by one point &
    BCO (boundary weakening), MBO (off-by-one) \\
    RF14 &
    Documentation--logic divergence &
    Comments or docstrings describe a scenario condition inconsistent with
    the implemented guard &
    SNO, ORO \\
    \bottomrule
    \end{tabularx}%
  }
\end{table}

All seven patterns except RF14 belong to $\mathcal{P}^H$ and therefore block
acceptance when detected.  RF14 is recorded as a medium-severity traceability
warning.  In the prevention evaluation, \emph{spec-confusion} trials inject
faults at the specification layer (operators SNO, ORO, MCO, BCO) and measure
whether a formally correct reference implementation is rejected against the
mutated specification---testing the formal checker's sensitivity to
specification drift rather than the pattern guard directly.  Nevertheless, the
patterns above are the static counterparts of the dynamic failures that
spec-confusion mutants induce: an implementation synthesised from a confused
reading of the specification will typically trigger one or more of RF01, RF02,
RF05, RF06, RF09, or RF11 before formal checking completes.

The precedence-sensitive guard structure of the hard benchmark
(Appendix~\ref{app:benchmark}) is deliberately constructed so that
specification-confusion faults manifest on overlap inputs: an implementation
that evaluates guards in reverse order may pass most randomly sampled inputs
yet fail on boundary witnesses generated by the SMT solver.  Patterns RF09 and
RF02 provide a complementary static signal when the confusion reduces branch
count or eliminates the default handler entirely.

% ---------------------------------------------------------------------------
\section{Implementation-Screening Patterns}
\label{sec:pat:is}
% ---------------------------------------------------------------------------

\emph{Implementation screening} targets defects introduced during code
generation that may not reflect a misunderstanding of the specification text
but nevertheless produce incorrect or structurally pathological behaviour.
Such faults include hard-coded outputs, missing return fields, neglected
external variables, and vacuous branch bodies.  They are evaluated in the
prevention protocol by injecting faults directly into reference implementations
(operators ICO, WRO, MBO, DRO) and measuring whether the hybrid acceptance
predicate rejects the mutated candidate.

Seven patterns serve as implementation-screening rules; four of them belong to
$\mathcal{P}^H$.  Table~\ref{tab:pat:impl-screening} summarises the mapping.

\begin{table}[htbp]
  \caption{Patterns associated with implementation-screening faults.}
  \label{tab:pat:impl-screening}
  \centering
  \thesistable{%
    \footnotesize
    \setlength{\tabcolsep}{3pt}
    \begin{tabularx}{\linewidth}{@{}p{0.06\linewidth}p{0.20\linewidth}X X@{}}
      \toprule
      \textbf{ID} & \textbf{Screening focus} & \textbf{Typical symptom} &
      \textbf{Related impl mutants} \\
    \midrule
    RF03 &
    Literal encoding &
    Magic numbers replace spec-derived expressions in return statements &
    WRO (wrong return constant) \\
    RF04 &
    Output completeness &
    Required output keys absent from the return dictionary &
    WRO (wrong return constant) \\
    RF07 &
    Spurious completeness &
    Fixed success value returned without conditional evaluation &
    WRO (wrong return constant) \\
    RF08 &
    Type discipline &
    Floating-point or untyped values returned where integers are required &
    --- \\
    RF10 &
    SOFL externality &
    Declared external variables never read or written &
    --- \\
    RF12 &
    Redundant assignment &
    Duplicate writes to the same output within one branch &
    --- \\
    RF13 &
    Vacuous branch &
    \code{else: pass} leaves a scenario handler without output assignment &
    DRO (else removal) \\
    \bottomrule
    \end{tabularx}%
  }
\end{table}

RF07 is the sole \emph{critical} pattern and represents the most severe
implementation-screening failure: a candidate that unconditionally returns a
fixed output passes any finite test suite whose witnesses happen to lie in
regions where that output is correct, yet violates the specification globally.
The pattern guard detects this pathology statically by requiring at least one
conditional branch before any constant success return.

Patterns RF03, RF08, RF10, and RF12 are rated medium or low severity and do
not participate in $\mathcal{P}^H$.  They are included in the catalogue because
empirical analysis of LLM-generated Agile-SOFL implementations revealed high
incidence of these coding defects, and their presence in the audit trace
provides actionable feedback for the repair loop even when acceptance is not
blocked.  For example, RF04 and RF07 together cover the two dominant failure
modes of the WRO mutant (wrong primary output and structurally incomplete
returns), while RF13 and RF02 jointly detect the structural consequences of
DRO (removed else branch).

In operational terms, the pattern guard executes only after a candidate
achieves strict formal conformance ($\mathrm{Conf}(c,t) = 1$), ensuring that
implementation-screening adds negligible overhead to definitively failing
candidates while providing a necessary second line of defence against
pathological implementations that satisfy all SMT-generated witnesses.
